% Fichier généré automatiquement par tools/generate_corrections.py.
% Source : SLCI/SLCI-03-SchemaBlocs/39_SeineMusicale/39_SeineMusicale.tex
% Statut : complet (3/3 question(s) renseignée(s), 0 reprise(s)).
\begin{corrigebox}[Corrigé extrait de la banque]
\CorrigeQuestion{1}
Réduction de la boucle du moteur à courant continu : 
$\dfrac{\Omega_m(p)}{U_m(p)}=\dfrac{\dfrac{k_c}{R+Lp}\dfrac{1}{J_{eq}p}}{1+\dfrac{k_c}{R+Lp}\dfrac{k_e}{J_{eq}p}}$
$=\dfrac{k_c}{\left(R+Lp\right)J_{eq}p+k_ek_c}$.

On a alors, 

$
\dfrac{X_{ch}(p)}{\Omega_c(p)} =
K_a  \dfrac{CK_h \dfrac{k_c}{\left(R+Lp\right)J_{eq}p+k_ek_c}}{1+CK_hK_{\text{capt}} \dfrac{k_c}{\left(R+Lp\right)J_{eq}p+k_ek_c}}
$ 

$
=
K_a \dfrac{CK_h k_c}{\left(R+Lp\right)J_{eq}p+k_ek_c+CK_hK_{\text{capt}} k_c}
$

$
=
\dfrac{K_a }{\left( k_ek_c+CK_hK_{\text{capt}} k_c\right)} \dfrac{CK_h k_c}{\dfrac{J_{eq}\left(R+Lp\right)}{k_ek_c+CK_hK_{\text{capt}} k_c}p+1}
$.
\CorrigeQuestion{2}
Par lecture directe du schéma-blocs, on a 
$\Omega_m(p) = \dfrac{1}{J_{eq}p}\left(C_{\text{pert}}(p) + C_m(p)\right)$.

De plus, 
$C_m(p) = \left(U_m(p)-k_e\Omega_m(p)\right) \dfrac{k_c}{R+Lp}$
 et $U_m(p)=\varepsilon(p) C K_h = -\Omega_m(p) C K_h K_{\text{capt}}$.
 
 On a donc, 
 
 $\Omega_m(p) = \dfrac{1}{J_{eq}p} C_{\text{pert}}(p) 
 + \dfrac{1}{J_{eq}p} \left(-\Omega_m(p) C K_h K_{\text{capt}}-k_e\Omega_m(p)\right) \dfrac{k_c}{R+Lp}$.

$\Leftrightarrow 
\Omega_m(p) = \dfrac{1}{J_{eq}p} C_{\text{pert}}(p) 
 + \dfrac{1}{J_{eq}p} \Omega_m(p) \left(- C K_h K_{\text{capt}}-k_e\right) \dfrac{k_c}{R+Lp}$
 
 $\Leftrightarrow 
\Omega_m(p)\left(1+\dfrac{1}{J_{eq}p} \left( C K_h K_{\text{capt}}+k_e\right) \dfrac{k_c}{R+Lp}\right)= \dfrac{1}{J_{eq}p} C_{\text{pert}}(p) 
 $

$\Leftrightarrow 
\dfrac{\Omega_m(p)}{C_{\text{pert}}(p)}
 = \dfrac{\dfrac{1}{J_{eq}p}}{\left(1+\dfrac{1}{J_{eq}p} \left( C K_h K_{\text{capt}}+k_e\right) \dfrac{k_c}{R+Lp}\right)}
 $

$\Leftrightarrow 
\dfrac{\Omega_m(p)}{C_{\text{pert}}(p)}
 = \dfrac{R+Lp}{J_{eq}p\left(R+Lp\right)+ \left( C K_h K_{\text{capt}}+k_e\right) k_c}
 $

$\Leftrightarrow 
\dfrac{\Omega_m(p)}{C_{\text{pert}}(p)}
 = \dfrac{R}{\left( C K_h K_{\text{capt}}+k_e\right) k_c} \dfrac{1+\dfrac{L}{R}p}{\dfrac{J_{eq}}{\left( C K_h K_{\text{capt}}+k_e\right) k_c}p\left(R+Lp\right)+ 1}
 $.
 
 Par identification, on a alors : 
 $\alpha = - \dfrac{R}{\left( C K_h K_{\text{capt}}+k_e\right) k_c}$, 

 $\tau = \dfrac{L}{R}$

$\gamma = \dfrac{R J_{eq}}{\left( C K_h K_{\text{capt}}+k_e\right) k_c} $

$\delta = \dfrac{LJ_{eq}}{\left( C K_h K_{\text{capt}}+k_e\right) k_c}$.
\CorrigeQuestion{3}
D'une part, $\Omega_m(p) = H_f(p) \Omega_c(p)$ quand il n'y a pas de perturbation.
D'autre part, $\Omega_m(p) = H_r(p) C_{\text{pert}}(p)$ quand il n'y a pas de perturbation.

Par superposition, on a donc $\Omega_m(p) = H_f(p) \Omega_c(p) + H_r(p) C_{\text{pert}}(p)$.

Par suite, $X_{ch}(p)=\left(H_f(p) \Omega_c(p) + H_r(p) C_{\text{pert}}(p)\right) \dfrac{DK_{\text{red}}}{2p}$.

\begin{enumerate}
\item $ H_f(p)=
\dfrac{K_a }{\left( k_ek_c+CK_hK_{\text{capt}} k_c\right)} \dfrac{CK_h k_c}{\dfrac{J_{eq}\left(R+Lp\right)}{k_ek_c+CK_hK_{\text{capt}} k_c}p+1} $.
\item  $\alpha = - \dfrac{R}{\left( C K_h K_{\text{capt}}+k_e\right) k_c}$,  $\tau = \dfrac{L}{R}$, 
$\gamma = \dfrac{R J_{eq}}{\left( C K_h K_{\text{capt}}+k_e\right) k_c} $, 
$\delta = \dfrac{LJ_{eq}}{\left( C K_h K_{\text{capt}}+k_e\right) k_c}$.
\item $X_{ch}(p)=\left(H_f(p) \Omega_c(p) + H_r(p) C_{\text{pert}}(p)\right) \dfrac{DK_{\text{red}}}{2p}$.
\end{enumerate}
\end{corrigebox}
